\ficha{Criação, 1906: os vencidos}

\bloco{Fontes lidas}

\textit{Documentos Parlamentares}, volume de 1906, lidos na íntegra nos trechos
indicados: voto em separado de Serzedelo Corrêa, de 15 de agosto (p.~72-79);
seus discursos de 30 de agosto (p.~212-226), 13 de setembro (p.~310-322) e 26 de
setembro (p.~448-455) e a declaração de voto de 29 de novembro (p.~611-612);
voto em separado de Paula Ramos (p.~80-99); discurso de Barbosa Lima de 26 de
setembro (p.~455-467); discursos de Alcindo Guanabara de 28 de agosto
(p.~184-211) e 25 de setembro (p.~440-447). As votações nominais estão nas
p.~345 (14 de setembro, 118 a 17), 544 (8 de outubro, 115 a 25) e 602 (Senado,
26 de novembro, 27 a 5).

\bloco{Síntese}

Os que votaram contra a lei de 1906 não formam um bloco doutrinário. Li quatro
posições distintas. Barbosa Lima e Paula Ramos querem o par de 1846, alcançado
por economias, resgate e fundo de garantia. Serzedelo Corrêa aceita a quebra do
padrão e até uma caixa de emissão, mas julga que a hora não chegou. Alcindo
Guanabara quer a fixação e a quer mais radical que a do projeto, com quebra
definitiva a 12 dinheiros; vota a favor do princípio e só se volta contra o
projeto depois das emendas de 24 de setembro. O que une esses nomes no
resultado da votação não é uma tese sobre a moeda.

\bloco{Serzedelo Corrêa: objeção de oportunidade}

O voto em separado abre declarando que \tr[P:1:78]{A sua divergencia é profunda
e completa} \dpa{p.~72} e sustenta que \tr[P:1:79]{não se póde fixar o preço de
mercadoria alguma só por simples vontade dos poderes publicos ou por effeito de
lei} \dpa{p.~73}. A conclusão é que \tr[P:1:85]{basta o funccionamento do fundo
de resgate para dar-nos a estabilização} \dpa{p.~79}.

A leitura integral dos discursos mostra que a objeção à taxa é de momento, não
de nível. Em 30 de agosto ele descreve três caminhos para abolir o curso
forçado e conclui que nenhum é praticável naquela hora, por causa da política
de empréstimos do quadriênio que terminava:

\cit[P:1:228]{O governo actual preferiu uma politica de lantejoulas a uma
politica de economias}{\dpa{p.~222}}

Sobre o aparelho do projeto, que não agiria na baixa do câmbio:

\cit[P:1:229]{a caixa é um organismo castrado, nada produz, não
funcciona}{\dpa{p.~223}}

Sobre a taxa, a de 15 lhe parece alta demais para se sustentar:

\cit[P:1:232]{a taxa de 15 d. me parece instavel, perigosa, incapaz de
manter-se}{\dpa{p.~226}}

e a de 12, que seria segura, \tr[P:1:319]{é depressiva, é nesta hora uma
expoliação, é baixa de mais} \dpa{p.~313}. Em 13 de setembro ele expõe o
próprio programa. Não é contra uma caixa de emissão:

\cit[P:1:320]{Acho que o processo do Japão é o verdadeiro: ao lado do resgate, a
fundação de um banco ou caixa de emissão de moeda conversivel}{\dpa{p.~314}}

A condição é que a emissão conversível substitua papel inconversível:
\tr[P:1:324]{Faça-se isso e a caixa é um apparelho complementar do mecanismo de
1899. Como está, nunca.} \dpa{p.~318}. A quebra do padrão deveria ser
definitiva, feita em ocasião oportuna e \tr[P:1:327]{a uma taxa sufficientemente
alta de modo a não causar prejuizos} \dpa{p.~321}.

Na declaração de 29 de novembro, depois das emendas da Câmara e do Senado, ele
reconhece que o projeto \tr[P:1:617]{é a expressão exacta do mecanismo que
formulei em meu voto vencido} \dpa{p.~611} e mantém o voto contrário por outra
razão: funcionando a Caixa, \tr[P:1:617]{a quebra do padrão a 15 se deve
considerar definitiva}, e \tr[P:1:617]{a taxa de 15 é evidentemente depressiva}
\dpa{p.~611}. As duas avaliações são compatíveis: 15 é taxa alta demais para
se manter em 1906 e baixa demais para valer como padrão permanente. O caminho
que ele propõe inclui a \tr[P:1:618]{valorização da que temos, especialmente o
café} \dpa{p.~612}. Serzedelo, vencido na Caixa, não é adversário da
valorização do café.

No mesmo discurso de 13 de setembro ele classifica os adversários do projeto.
Paula Ramos e Barbosa Lima \tr[P:1:321]{querem a politica rigorosa do
quatriennio passado, querem a paridade da lei de 1846}; Alcindo Guanabara
\tr[P:1:321]{quer a caixa de conversão e quer com quebra do padrão a 12}
\dpa{p.~315}. E diz que os que se identificaram com a política de Murtinho e
Campos Sales não podem aceitar o mecanismo: \tr[P:1:319]{esses nunca, salvo por
uma verdadeira evolução politica} \dpa{p.~313}.

\bloco{Barbosa Lima e Paula Ramos: o par de 1846}

Paula Ramos resume a recusa no voto em separado: \tr[P:1:105]{A caixa de que
trata o projecto não fixa e nem converte} \dpa{p.~99}. Barbosa Lima, em 26 de
setembro, leva o argumento da estabilidade ao absurdo:

\cit[P:1:462]{Estavel a 24? Estavel a 15? Estavel a 12? Estavel a 9? Ou estavel
a 3?}{\dpa{p.~456}}

Ele concede que o mecanismo quebra o padrão de fato, tese que Campista negava:
\tr[P:1:464]{Os effeitos praticos são de quebra do padrão, de mudança do par.
Nós quebramos o padrão} \dpa{p.~458}. Atribui aos partidários da fixação a
convicção de que \tr[P:1:465]{aquillo de que este paiz mais precisa, agora,
depois do quatriennio Campos Salles, é cambio baixo} \dpa{p.~459}. Invoca
Joaquim Murtinho, para quem fixar o câmbio para converter seria
\tr[P:1:470]{uma deshonestidade por parte do Estado} \dpa{p.~464}, e fecha com
a fórmula: \tr[P:1:472]{Eu, que não sou da colligação, sustento o programma da
colligação} \dpa{p.~466}.

A mesma citação de Murtinho já aparecia no resumo que a \textit{Gazeta de
Notícias} fez do discurso de 24 de agosto (\jor{gn19060825}{Gazeta de
Notícias}{25 ago. 1906}{3}), sessão que os anais trazem só em discurso
indireto. É argumento recorrente de Barbosa Lima, não divergência de data entre
as fontes.

\bloco{Alcindo Guanabara: a favor do princípio, contra o aparelho}

O discurso de 28 de agosto conclui a primeira parte afirmando a necessidade da
fixação:

\cit[P:1:204]{o projecto em debate, cujo intuito é conseguir este desideratum,
não é só conveniente, como necessario}{\dpa{p.~198}}

O que ele nega é a eficácia do aparelho sem quebra do padrão. Calcula que, no
ritmo de 1900 a 1904, o resgate levaria \tr[P:1:199]{cerca de 68 annos}
\dpa{p.~193}, e propõe a quebra definitiva a 12 dinheiros, porque
\tr[P:1:217]{a média geral desses 14 annos é de 11 9/32} \dpa{p.~211}.

Em 25 de setembro ele se descreve em posição singular entre os impugnadores,
\tr[P:1:447]{acceitando o principio que nelle se consubstancia} \dpa{p.~441},
e volta à tribuna por causa das emendas que autorizam a Caixa a operar em
câmbio com o fundo de garantia: \tr[P:1:452]{Entregar o fundo de garantia á
Caixa de Conversão para operar em cambio é sem duvida pôl-o em risco}
\dpa{p.~446}. Daí a frase de fecho: \tr[P:1:453]{o que era inutil passou a ser
perigoso para o paiz} \dpa{p.~447}.

\bloco{Achados}

\begin{enumerate}[leftmargin=*,nosep]
\item O voto \textit{sim} de Alcindo em 14 de setembro, que conferi na imagem
da p.~345, não contradiz seus discursos. Ele aceita o princípio e vota o
art.~1º; a virada é posterior e tem causa datada, as emendas de 24 de setembro.
Na votação de 8 de outubro ele não aparece.
\item A frase de Serzedelo sobre a taxa \textit{depressiva}, de novembro, não
pode ser lida isoladamente. Em agosto e setembro o risco que ele aponta é o de
a taxa de 15 não se sustentar.
\item A oposição de 1906 reúne valorizadores do par (Barbosa Lima, Paula
Ramos), um defensor da quebra definitiva em outra data e a outra taxa
(Serzedelo) e um estabilizador a taxa mais baixa que a do projeto (Alcindo).
Uma escala única de favorável ou contrário à Caixa apagaria essas diferenças.
\end{enumerate}

\bloco{Limites}

O discurso de Barbosa Lima de 26 de setembro traz nos anais a nota de que não
foi revisto pelo orador. O de 24 de agosto só existe em resumo indireto. As
páginas 226 e 313 do volume têm OCR ruidoso mesmo na camada corrigida.
